% A137: full source audited by root, local_preprint and birational_main.
\subsection{An explicit completed tripod pair excluded by a finite geometric quotient}

\begin{statement}{A.137}{Proved}
{The harmonic Brauer obstruction on the actual tripod}
Put
\[
 R=\frac{529}{382},\qquad h=\frac12,\qquad
 s=155564491,
 \qquad (a_0,a_1,a_h)=\bigl(R,R-1,(R-h)s^2\bigr).
\tag{A.137.1}
\]
Let \((A_0,A_1,A_h)\) be the completed harmonic triple
constructed in A.134, with its fixed ordinary finite
centres \(T_p\), and define
\[
 A=-4A_0A_1,\qquad B=4A_h^2,\qquad
 x_p=4T_p(1-T_p).
\tag{A.137.2}
\]
This is an ordinary completed unit-equation pair on
\(U=\mathbf P^1_\Q\setminus\{0,1,\infty\}\).
Both coordinates have infinite support.  It has no
full Selmer section with those complete Kummer coordinates.

The obstruction occurs in the following actual finite
geometric quotient.  Use the degree-two cover
\(V=\mathbf P^1_\Q\setminus\{0,1,h,\infty\}\to U\)
of A.107, compatible rational base points, and
\[
 \begin{gathered}
 \Lambda=\pi_1(V_{\overline\Q})\triangleleft
       \Pi=\pi_1(U_{\overline\Q}),\qquad
 H=\ker(\Lambda\longrightarrow\Lambda^{\rm ab}/4),\\
 K=\overline{H^2[H,H]},\qquad Q=\Pi/K,\qquad
 |\Lambda/K|=2^{135},\quad |Q|=2^{136}.
 \end{gathered}
\tag{A.137.3}
\]
There is no continuous section of \(E_U/K\to G_\Q\)
whose Kummer class of \(1-t\) modulo eight is
\([4a_h^2]\), whose localization at every finite prime
is geometrically conjugate to the image of the ordinary
point \(x_p\), and whose real localization is the image
of an ordinary or cuspidal real section.
All these classes and local images are defined on the
displayed quotient.
The actual arithmetic maps
\[
 E_U/\Pi'''\longrightarrow E_U/K,\qquad
 E_U/\Gamma_{18}\Pi\longrightarrow E_U/K
\tag{A.137.4}
\]
therefore give the same exclusion for locally matching
sections of the third-derived quotient and of the
lower-central quotient of class17.  The geometric group
$Q$ has exact nilpotency class17, with a specified
nonzero weight-seventeen commutator below.
This assertion concerns this particular pair; it does
not assert that every infinite-support pair fails at
that quotient.
\end{statement}

\begin{proof}
We retain the nonsplit arithmetic group of the
quartic cover throughout the proof.
The actual cover and its Brauer class from A.134 are
\[
 \begin{gathered}
 C^\circ:\quad T-i=a_iX_i^4\quad(i=0,1,h),\qquad X_i\ne0,\\
 c=-\frac{338}{191},\qquad d=7\cdot17s=18512174429,\qquad
 \alpha=(T-c,d)\in\operatorname{Br}(\overline C)[2].
 \end{gathered}
\tag{A.137.5}
\]
Here \(\overline C\) is the smooth proper curve, and
\(C^\circ\) removes precisely the fibres over
\(0,1,h,\infty\).  In particular the fibre over \(c\)
is retained.  A.134 proves that the class is unramified
on this proper curve and that its local evaluation is
constant at every place.  The sum of these constants is
\(1/2\), with its only nonzero value at the dyadic place.
Its real value is zero.  The same holds at real boundary
points because the class extends to \(\overline C\).

\emph{A finite descent of the Brauer class without an arithmetic splitting.}
Choose a lift \(\widetilde\alpha\in
H^2_{\rm et}(\overline C,\mu_2)\) through the Kummer
sequence, and restrict it to \(C^\circ\).
The comparison map gives
\[
 H^2(\pi_1(C^\circ),\F_2)
       \xrightarrow{\ \sim\ }
 H^2_{\rm et}(C^\circ,\mu_2).
\tag{A.137.6}
\]
Its precise justification is the one in (A.93.7),
which does not require the arithmetic extension to
split: the geometric affine curve has free profinite
fundamental group \(H\), its finite-coefficient etale
cohomology vanishes above degree one, and the etale
and group Hochschild--Serre spectral sequences have
the same geometric terms and comparison maps.
The geometric statements follow from the finite-cover
and etale--Betti comparisons and the curve localization
sequence cited and proved there.
Let \(E=\pi_1(C^\circ)\), and let
\(e\in H^2(E,\F_2)\) correspond to this restriction.

Represent \(e\) by its continuous central extension
\[
               1\longrightarrow\F_2\longrightarrow
                  \widetilde E\longrightarrow E\longrightarrow1.
\tag{A.137.7}
\]
The inverse image of \(H\) splits over \(H\), since
a free profinite group is projective.  Explicitly,
choose lifts of its finite free generating set in
that inverse image; the defining universal property
extends them to a continuous homomorphism
\(j:H\to\widetilde E\) splitting the projection.
For \(z\in\widetilde E\), let \(\tau_z\) be conjugation
by its image on \(H\).  Both maps
\[
 h'\longmapsto zj(h')z^{-1},\qquad
 h'\longmapsto j(\tau_z(h'))
\tag{A.137.8}
\]
are continuous homomorphisms lifting \(\tau_z\).
Their quotient lies in the central group \(\F_2\),
so it is a continuous homomorphism \(H\to\F_2\).
Every such homomorphism kills \(K\).
Thus \(j(K)\) is a closed normal subgroup of
\(\widetilde E\), and it intersects the central
\(\F_2\) trivially.
Consequently
\[
 \begin{gathered}
 1\longrightarrow\F_2\longrightarrow
       \widetilde E/j(K)\longrightarrow E/K\longrightarrow1,\\
 \widetilde E\xrightarrow{\ \sim\ }
       E\mathop{\times}_{E/K}\bigl(\widetilde E/j(K)\bigr),
       \qquad z\longmapsto(\overline z,[z])
 \end{gathered}
\tag{A.137.9}
\]
are an actual central extension and a topological
isomorphism.  Injectivity of the second map follows
from \(j(K)\cap\F_2=1\).  For surjectivity, correct a
representative of its second coordinate by an element
\(j(k)\), \(k\in K\), to give the prescribed first
coordinate.  Compactness and the Hausdorff property
then give the claimed topology.
The class \(e_Q\in H^2(E/K,\F_2)\) of the first
extension inflates to \(e\).
No section of \(E\to G_\Q\) was used.

\emph{The quartic cover is forced by finite tripod data.}
We first prove a rational fourth-power detection fact.
Let \(b\in\Q^\times\) be a fourth power in \(\Q_p\)
for all sufficiently large primes in each of the
progressions \(p\equiv5\pmod8\) and
\(p\equiv3\pmod8\).
Set
\[
                L=\Q(\zeta_8),\qquad F=L(b^{1/4}).
\tag{A.137.10}
\]
This is a finite Galois extension of \(\Q\).
Every element of \(\operatorname{Gal}(F/\Q)\) whose
restriction to \(L\) is \(\sigma_5:\zeta_8\mapsto\zeta_8^5\)
is a Frobenius element, up to conjugacy, at infinitely
many unramified primes in the first progression.
At those primes the fourth-root equation has a root,
so this element fixes one of the four geometric roots.
It fixes \(\mu_4\), and hence its action on the root
set is a translation.  A translation fixing one root
fixes every root.
This is true for every lift of \(\sigma_5\).
Taking quotients of two such lifts makes every element
of \(\operatorname{Gal}(F/L)\) fix every root, so
\(F=L\).  Moreover all four roots are fixed by
\(\sigma_5\), and therefore lie in its actual fixed
field \(\Q(i)\).

Write a root as \(u+vi\), with \(u,v\in\Q\).
The imaginary part of its fourth power is
\(4uv(u^2-v^2)\).
Its vanishing gives either \(b=w^4\) or
\(b=-4w^4\), with \(w\in\Q^\times\).
The latter has no square root at any sufficiently
large prime congruent to three modulo four, since
its squareclass is \(-1\).
The second progression excludes that alternative.
This proves \(b\in(\Q^\times)^4\) with the exact
field and root-set maps just used.

Now suppose a section as specified in the statement
is given.  The mod-two character of \(1-t\) is zero,
since its mod-eight class is \([4a_h^2]\).
The actual degree-two cover therefore lifts the
section to \(E_V/K\), as in (A.107.2).
This quotient retains each old Kummer coordinate
modulo eight: its geometric character is zero on
\(K\), because a character of \(\Lambda\) takes
\(H\) into \(4\widehat\Z\), and takes
\(H^2[H,H]\) into \(8\widehat\Z\).
Also \(1-t=4(T-h)^2\), so its character on \(K\)
is divisible by sixteen.  Thus every finite Kummer
class used here is defined on the indicated quotient.

Write \(C_0,C_1,C_h\in\Q^\times/(\Q^\times)^8\)
for the lifted section's old coordinates.
The literal coordinate identity gives
\[
                    C_h^2=[a_h^2]
                       \quad\text{in }\Q^\times/(\Q^\times)^8.
\tag{A.137.11}
\]
The rational valuation and sign decomposition makes
\(C_h=[a_h]\) or \([-a_h]\) modulo four: the valuation
of a quotient whose square is an eighth power is
divisible by four at every rational prime, and its
only remaining factor is a sign.
Apply the single global deck involution if needed,
so that \(C_h=[a_h]\) modulo four.

At each finite prime the local lift is geometrically
conjugate to the image of one of the two actual points
\(T_p,1-T_p\) above \(x_p\).
Indeed the fixed sheets of the degree-two finite etale
cover are precisely those two local points; conjugacy
in its two cosets gives exactly these two possibilities.
By A.134, \([T_p-h]=[a_h]\) modulo four locally.
At an odd prime \(p\not\equiv1\pmod8\), the element
\(-1\) is not a fourth power in \(\Q_p\).
Hence the other point, whose \(T-h\) value is
\(-(T_p-h)\), cannot have the chosen class \(C_h\).
The lifted point is therefore \(T_p\) at every such
prime.  In particular
\[
                 (C_0)_p=[a_0],\qquad
                 (C_1)_p=[a_1]
                \quad\bmod (\Q_p^\times)^4
\tag{A.137.12}
\]
in both progressions used in (A.137.10).
Choose rational representatives of \(C_i/a_i\).
The proved fourth-power detection fact applies to
each representative, so the global mod-four triple
is exactly \(([a_0],[a_1],[a_h])\).

The mod-four root torsor for the equations of
\(C^\circ\) consequently becomes trivial on the
lifted section.  In the finite-cover correspondence,
this is the assertion that the section fixes a
geometric sheet, or equivalently that, after geometric
conjugation, its image lies in the actual subgroup
\(E/K\subset E_V/K\).
The subgroup statement is literal because \(K\)
is contained in the geometric subgroup \(H\) of
the cover.  We have obtained a section
\(s_C:G_\Q\to E/K\).

\emph{The global contradiction and its local comparisons.}
Pull back (A.137.9) along \(s_C\).
The resulting class
\[
                    s_C^*e_Q\in H^2(\Q,\mu_2)
                                      =\operatorname{Br}(\Q)[2]
\tag{A.137.13}
\]
has the following exact local evaluations.
At a finite prime, its local section is the image
of a point of \(C^\circ(\Q_p)\).
To see this, lift its finite geometric conjugating
element to the geometric fundamental group and
use the fixed sheet of the actual mod-four cover;
the resulting Galois-fixed sheet is an actual local
point in that fibre.
The cohomology comparison and the fibre-product
isomorphism in (A.137.9) then identify the restriction
of (A.137.13) with \(\alpha(P_p)\).
The evaluation is independent of the selected root
sheet by A.134's constant-evaluation theorem.

At the real place the same comparison applies to
an ordinary lift.  For a cuspidal lift, the finite
cover's inertia orbits give a real point on the
proper normalization above that cusp.  Because
\(\widetilde\alpha\) was chosen on \(\overline C\),
its restriction along the tangential section is
the pullback at that proper real point: the map
to the proper curve kills the cusp inertia.
The induced Brauer class is again its value of
\(\alpha\), namely zero.
This proves the required comparison also for the
allowed real cuspidal images.

The local invariants of (A.137.13) therefore sum
to \(1/2\), by (A.137.5) and A.134.
Global reciprocity for the actual rational Brauer
class (A.137.13) makes their sum zero.
This contradiction proves the finite-quotient
nonexistence assertion.

\emph{The completed pair and its quotient consequences.}
The identities (A.137.2) give the completed pair
and its ordinary local centres by literal
localization, exactly as in A.107.
The fresh coefficient primes of coordinates zero
and one in A.134 give infinite support of \(A\);
those of coordinate \(h\) give twice their positive
valuations in \(B\).
At the real place its actual ordinary centre is
\(4R(1-R)=-77763/36481\).
Its mod-eight \(B\)-class is \([4a_h^2]\), since
the old completed coordinate has mod-four class
\([a_h]\).

A full Selmer section with this pair has the
ordinary local centres \(x_p\) at every finite
prime: the complete local Kummer map is injective,
and each of the three possible cusp labels is
excluded by the regular-coordinate identities in
the proof of (A.107.6).
Its real image has the stated Selmer property.
Projection to \(E_U/K\) would therefore give the
section just excluded.
This proves the full Selmer nonexistence.

For completeness, \(H\) is free of rank
\(1+64(3-1)=129\), giving \(|H/K|=2^{129}\)
and the orders in (A.137.3).
The characteristic definitions make \(K\)
normal under the complete arithmetic group.
Since \(\Pi'\subset\Lambda\),
\(\Pi''\subset\Lambda'\subset H\), and
\(\Pi'''\subset[H,H]\subset K\), the first
map in (A.137.4) is an actual quotient map.
The exact finite calculation below proves
$\Gamma_{18}Q=1$ and $\Gamma_{17}Q\ne1$, giving
the second map in (A.137.4).
The local images and the retained mod-eight
Kummer character pass through both maps.
This proves all stated quotient consequences.

% Insert within A.137 after its finite geometric quotient construction.
% This fragment proves the replacement Gamma_136 -> Gamma_18, and exact class17.
\paragraph{The exact nilpotency class of the actual finite quotient.}
The geometric group $Q=\Pi/K$ above has nilpotency
class exactly $17$.  Thus its actual lower-central
quotient map can be sharpened to
\[
                            \Pi/\Gamma_{18}\Pi\twoheadrightarrow Q.
\]
Here is the covering-graph calculation, including a
specific nonzero commutator of weight seventeen.

Choose the free generators $x,y$ of $\Pi$ so that
the index-two character has values $0,1$ on them.
The Schreier generators of $\Lambda$ are
$a=x$, $b=yxy^{-1}$ and $z=y^2$.
After passing to $\Lambda^{\rm ab}/4$ they commute
and have order four.
Conjugation by $y$ exchanges $a,b$ and fixes $z$.
Consequently the exact group $P=\Pi/H$ is
\[
 \begin{gathered}
 P=(C_4\times C_4)\rtimes C_8,\qquad |P|=128,\\
 (a,b,k)(a',b',k')
   =\bigl((a,b)+s^k(a',b'),\,k+k'\bigr),\\
       s(a',b')=(b',a'),\quad
       a,b,a',b'\in\Z/4,\quad k,k'\in\Z/8,\\
                       x=(1,0,0),\qquad y=(0,0,1).
 \end{gathered}
\]
This presentation has the actual subgroup
$\langle a,b,y^2\rangle=C_4^3$ and the indicated
index-two quotient, so its kernel is exactly $H$.
With $[u,v]=uvu^{-1}v^{-1}$ one has
$[x,y]=ab^{-1}$,
$[[x,y],y]=(ab^{-1})^2\ne1$,
and every weight-four commutator in $P$ is trivial.
Thus $P$ has class three.

Take the covering graph of the two-loop rose with
vertices $g\in P$ and directed edges
$e_{g,x}:g\to gx$ and $e_{g,y}:g\to gy$.
Over $\F_2$ set
\[
 \begin{gathered}
 C_1=\bigoplus_{g\in P}\bigl(\F_2e_{g,x}\oplus\F_2e_{g,y}\bigr),
       \qquad C_0=\bigoplus_{g\in P}\F_2\delta_g,\\
 \partial e_{g,x}=\delta_g+\delta_{gx},\qquad
 \partial e_{g,y}=\delta_g+\delta_{gy},\qquad
             h\cdot e_{g,u}=e_{hg,u}.
 \end{gathered}
\]
Left translation commutes with the boundary map.
The following is an explicit finite group:
\[
 \begin{gathered}
 \mathcal Q=\{(v,g)\in C_1\times P:
                          \partial v=\delta_g+\delta_1\},\\
 (v,g)(w,h)=(v+g\cdot w,gh),\qquad
                  (v,g)^{-1}=(g^{-1}\cdot v,g^{-1}),\\
                   X=(e_{1,x},x),\qquad Y=(e_{1,y},y).
 \end{gathered}
\]
The free-generator map $\Pi\to\mathcal Q$ is
the actual path map modulo two.
For an element of $H$ it records a closed path
in this graph.
A spanning tree gives $256-128+1=129$ independent
non-tree-edge cycles; the corresponding Schreier
loops form the free generators of $H$.
Hence its restriction induces the exact isomorphism
\[
        H/\overline{H^2[H,H]}\ \xrightarrow{\ \sim\ }
                       Z_1:=\ker(\partial:C_1\to C_0).
\]
For each endpoint $g$, choose a word reaching it.
Adding any cycle by an element of $H$ gives every
pair in the displayed fibre over $g$.
Thus $\Pi\to\mathcal Q$ is onto, its kernel is
exactly $K$, and it identifies $Q$ with this
specific group.
In particular $|Q|=128\cdot2^{129}=2^{136}$.
All these maps have finite targets and are
continuous for the profinite groups in A.137.

This gives the actual comparison with the established
relation-module construction.  In
\cite[Definition 2.1 and Theorem 3.7]{AgentDietrichHulpke2020},
take the finite group to be $P$, the free rank to be two,
and the prime to be two.  The map sending the path image
of a word $w$ to its image $\Psi_2(w)$ in that theorem
is a well-defined group isomorphism: the two kernels
are exactly $\overline{H^2[H,H]}$, as proved here and
in that theorem.  On the kernel, the edge maps
$e_{g,x}\mapsto(g,0)$ and $e_{g,y}\mapsto(0,g)$
identify our boundary sequence with the mod-two
relation sequence in
\cite[Proposition 2.8]{AgentYamin1991}; its displayed
boundary is right multiplication by $x-1,y-1$ on
the two left regular summands.
The specific lower-central calculation follows below.

Since $P$ has class three, $\Gamma_4Q$ is a
subspace of the elementary abelian cycle group
$Z_1$.
Its exact generating vectors are the $P$-translates
of the four actual commutators
\[
 [[[X,Y],X],X],\quad [[[X,Y],X],Y],\quad
 [[[X,Y],Y],X],\quad [[[X,Y],Y],Y].
\]
To justify completeness, the normal closure of
$[X,Y]$ is $\Gamma_2Q$.
If a normal subgroup is normally generated by
a list, its commutator with $\langle X,Y\rangle$
is normally generated by the commutators of
that list with $X,Y$: in the quotient killing
those latter commutators every original
generator is central.
Apply this twice.
On $Z_1$ conjugation depends only on $P$,
giving precisely the stated vector span.

For a $P$-stable subspace $W\subset Z_1$ the next
commutator subspace is exactly
\[
                  [W,Q]=(x-1)W+(y-1)W.
\]
Indeed $[(v,1),X]=(v+x\cdot v,1)$, and similarly
for $Y$.
The identity
$(gh-1)w=(g-1)(h\cdot w)+(h-1)w$
expresses the difference for every word in $x,y$
through these two images.
It also proves their sum is $P$-stable.

The retained integer-bitset elimination uses
the indexing
$g=a+4b+16k$ and $e_{g,x}=2g$, $e_{g,y}=2g+1$.
It starts with those four literal commutators
and their $128$ translates, and applies the
last displayed recurrence until zero.
Its exact dimensions are
\[
 \begin{gathered}
 \begin{array}{c|rrrrrrrr}
 j&4&5&6&7&8&9&10&11\\ \hline
 \dim_{\F_2}\Gamma_jQ&124&121&115&107&97&85&71&55
 \end{array}
 \\[6pt]
 \begin{array}{c|rrrrrrr}
 j&12&13&14&15&16&17&18\\ \hline
 \dim_{\F_2}\Gamma_jQ&41&29&19&11&5&2&0 .
 \end{array}
 \end{gathered}
\]
Each retained basis element is an unreduced
independent original generator rather than
an unspecified linear combination.
Its recorded provenance therefore remains
an actual commutator at every step.

Explicitly, start with $C_4=[[[X,Y],X],X]$
and successively replace $C$ by $[C,Z]$
for the thirteen letters
\[
                         Z:\quad \mathrm{XXYXYXYYYYYYY}.
\]
The resulting weight-seventeen element is
\[
                    \left(\sum_{g\in P}e_{g,x},\,1\right)\ne1.
\]
Its bitset is the $256$-bit integer with every
even-indexed bit equal to one and every
odd-indexed bit zero.
It is nonzero, is a cycle, and is fixed by
both left translations.
The final zero row proves $\Gamma_{18}Q=1$;
this explicit element proves $\Gamma_{17}Q\ne1$.
Together they prove the asserted exact class and
the sharpened quotient map.

The code
\path{code/harmonic_quotient_class_certificate.py},
its full output, and
\path{code/harmonic_quotient_class_data.json}
retain every basis vector and its commutator
provenance.
The computation also checks all $128^3$
associative products of $P$, every translated
edge boundary, every canonical word lift,
the four starting paths, and the equality of
the recorded deepest word with its displayed
cycle.  No enumeration of the $2^{136}$
elements of $Q$ is used.

\paragraph{An independent augmentation-ideal upper bound.}
There is a second exact verification of the upper bound.
Put $\mathcal R=\F_2[P]$ and let $I$ be its augmentation
ideal.  The independent program
\path{code/harmonic_quotient_augmentation_audit.py}
computes its powers in the regular module by
$v\mapsto v+xv$ and $v\mapsto v+yv$.
The identity for $gh-1$ above proves that their span
on $I^j$ is exactly $I^{j+1}$.
Its dimensions are
\[
 \begin{aligned}
 (\dim I^j)_{j=0}^{17}
   ={}&(128,127,125,121,115,107,97,85,71,\\
       &57,43,31,21,13,7,3,1,0).
 \end{aligned}
\]
The successive differences agree coefficient by
coefficient with $(1+t)^2(1+t^2)^3(1+t^4)^2$;
in particular the executed finite calculation proves
$I^{17}=0$ and $\dim I^{16}=1$ directly.

For a free word $w$, write its path image as
$(D(w),\nu(w))$.  The multiplication already proved gives
$D(uv)=D(u)+\nu(u)D(v)$.
With $g=\nu(u)$ and $h=\nu(v)$, the exact identities are
\[
 \begin{aligned}
 [g,h]-1
   &\;=\bigl((g-1)(h-1)-(h-1)(g-1)\bigr)g^{-1}h^{-1},\\
 D([u,v])
   &\;=(1-ghg^{-1})D(u)+(g-[g,h])D(v),\\
 1-ghg^{-1}&\;=-g(h-1)g^{-1},\qquad
 g-[g,h]=gh(1-g^{-1})h^{-1}.
 \end{aligned}
\]
They prove simultaneously by induction that
\[
 \nu(w)-1\in I^j,\qquad D(w)\in (I^{j-1})^2
                              \quad(w\in\Gamma_j\Pi).
\]
Indeed commutators of weights $a,b$ have endpoint
difference in $I^{a+b}$ and chain in $(I^{a+b-1})^2$.
Products and inverses preserve these ideals.
Normal conjugation also preserves them by the literal formula
\[
 D(vuv^{-1})
   =(1-\nu(v)\nu(u)\nu(v)^{-1})D(v)+\nu(v)D(u).
\]
The finite target makes the assertion valid on
the closed profinite lower-central subgroups.
For $w\in\Gamma_{18}\Pi$, its endpoint is one
because $P$ has class three, and its chain is zero
because $I^{17}=0$.
The exact path kernel therefore gives
$\Gamma_{18}\Pi\subset K$ independently of the
preceding lower-central row-space calculation.
The same independent program expands the specified
weight-seventeen word into $196606$ letters and
traces its edges one by one, obtaining all $128$
$x$-edges and no $y$-edges.
This separately verifies its nonzero image.
The exact class of this detecting quotient does
not assert that the constructed pair first fails
at class seventeen.

\end{proof}
